% Source: https://artofproblemsolving.com/community/c3h2410306_something_i_discovered
% Stable identity: aops-topic-2410306-q01
Give each unit cube coordinates $(x,y,z)$ with $x,y,z\in\{1,2,\dots,n\}$. A unit cube has a face on the left or right face of the large cube exactly when $x=1$ or $x=n$, and then exactly one such face, because $1\ne n$ when $n\ge2$; similarly for $y$ and $z$. Call a coordinate \emph{extreme} if it is $1$ or $n$. Then the number of faces of a unit cube on the surface equals the number of its extreme coordinates.

Each coordinate is extreme in $2$ ways and not extreme in $n-2$ ways. A unit cube with exactly $k$ extreme coordinates is chosen by picking which $k$ coordinates are extreme, $\dbinom3k$ ways, and their values, $2^k(n-2)^{3-k}$ ways:
\begin{itemize}
\item $k=3$: $2^3=8$ (the corners);
\item $k=2$: $3\cdot4(n-2)=12(n-2)$ (along the $12$ edges);
\item $k=1$: $3\cdot2(n-2)^2=6(n-2)^2$ (on the $6$ faces);
\item $k=0$: $(n-2)^3$ (the inside).
\end{itemize}
By the binomial theorem the total is $\sum_{k=0}^3\dbinom3k2^k(n-2)^{3-k}=(2+(n-2))^3=n^3$, as it must be. $\blacksquare$

\begin{asy}
usepackage("amsmath");
size(5.5cm);
defaultpen(fontsize(10pt));
int n=4;
pair pr(real x,real y,real z){return (x+0.5*y,z+0.36*y);}
pen c3=rgb(0.95,0.45,0.45), c2=rgb(1,0.8,0.45), c1=rgb(1,1,0.65);
int ext(int t){return (t==0 || t==n-1) ? 1 : 0;}
pen col(int k){return k==3 ? c3 : (k==2 ? c2 : c1);}
for(int i=0;i<n;++i) for(int k=0;k<n;++k){
  fill(pr(i,0,k)--pr(i+1,0,k)--pr(i+1,0,k+1)--pr(i,0,k+1)--cycle,col(1+ext(i)+ext(k)));
  draw(pr(i,0,k)--pr(i+1,0,k)--pr(i+1,0,k+1)--pr(i,0,k+1)--cycle,gray);}
for(int i=0;i<n;++i) for(int j=0;j<n;++j){
  fill(pr(i,j,n)--pr(i+1,j,n)--pr(i+1,j+1,n)--pr(i,j+1,n)--cycle,col(1+ext(i)+ext(j)));
  draw(pr(i,j,n)--pr(i+1,j,n)--pr(i+1,j+1,n)--pr(i,j+1,n)--cycle,gray);}
for(int j=0;j<n;++j) for(int k=0;k<n;++k){
  fill(pr(n,j,k)--pr(n,j+1,k)--pr(n,j+1,k+1)--pr(n,j,k+1)--cycle,col(1+ext(j)+ext(k)));
  draw(pr(n,j,k)--pr(n,j+1,k)--pr(n,j+1,k+1)--pr(n,j,k+1)--cycle,gray);}
draw(pr(0,0,0)--pr(n,0,0)--pr(n,0,n)--pr(0,0,n)--cycle,linewidth(0.9));
draw(pr(0,0,n)--pr(0,n,n)--pr(n,n,n)--pr(n,0,n),linewidth(0.9)); draw(pr(n,0,0)--pr(n,n,0)--pr(n,n,n),linewidth(0.9));
fill(box((7.2,4.2),(7.6,4.6)),c3); label("3 faces",(7.6,4.4),E,fontsize(8pt));
fill(box((7.2,3.4),(7.6,3.8)),c2); label("2 faces",(7.6,3.6),E,fontsize(8pt));
fill(box((7.2,2.6),(7.6,3.0)),c1); label("1 face",(7.6,2.8),E,fontsize(8pt));
\end{asy}

(The sketch shows the visible faces for $n=4$: $8$ corner cubes, $12\cdot2=24$ edge cubes, $6\cdot4=24$ face cubes and $8$ hidden inner cubes, $64$ in all.)
